\subsection{零 $b^2$ 列の完成と三固有値全体の排除}
\label{subsec:cubic-third-root-nonunit-zero-column}
\begin{theorem}[最後の零 $b^2$ 列も空]
\label{thm:cubic-third-root-nonunit-zero-column-empty}
非単位重複rootの $p=r=0$ の全atlasはH3を与えない。
\end{theorem}
\begin{proof}
この場合、重み1の関係は二つの純 $Uc+cU$ 関係である。
どちらかの射影に核があれば純 $uc$ 又は $cu$ の零積となるので、整域性から両射影は同型。
従って $c$ は $U$ 上の正常元であり、可逆線型写像 $F$ を用いて
$c u=F(u)c$ と書ける。
自己対称middleは $RJR^t$、$R=\binom{-F}{I_2}$、$J$ は対称可逆と書け、
元巡回式の二主blockの比例から
\[
 FJF^t=-\omega^2J
\]
となる。

$m=a_{10},n=a_{18},T=a_{21}$ とし、$L=mx+ny$ と置く。
元middleの最後二列は正確に
\[
 Lb=T c^2,\qquad bL=\omega^2T c^2
\tag{Z0}
\]
を与える。
$m=n=0$ は不可能である。$T\ne0$ なら純 $c^2$ 関係、
$T=0$ なら二次元の純 $U^2$ 関係pencilがあり、いずれも整域性に反する。
従って $L\ne0$、この二列への射影rankは2である。
$T=0$ なら $c^2$ 行全体が零となり、自己対称middleのcross blockも零になる。
しかし上のgraph表示ではcrossは $-FJ$ で可逆なので矛盾する。
従って $T\ne0$。

(Z0)の結合則から
$c^2L=\omega^2Lc^2$、従って
\[
 F^2(L)=\omega^2L
\]
となる。
$F(L)=\lambda L$ の場合は、$c^2c=cc^2$ と (Z0) により
$L(bc-\lambda cb)=0$、従って $bc=\lambda cb$ である。
$c$ は全 $B$ の正常元となり、正常正則一次商 $B/(c)$ は
独立な $L,b$ に対して $Lb=bL=0$ を持つ。
補題\ref{lem:cubic-normal-linear-opposite-products}のHilbert矛盾に反する。

従って $L,F(L)$ は独立である。
Cayley--Hamiltonを $L$ に作用させると
\[
 (\operatorname{tr}F)F(L)=(\det F+\omega^2)L
\]
なので $\operatorname{tr}F=0,\det F=-\omega^2$。
$F$ は相異なる固有値 $\omega,-\omega$ を持ち、$U$ の基底をその固有基底へ取れる。
この基底で $FJF^t=-\omega^2J$ は $J$ の対角成分を零にし、
非零off-diagonalだけが残る。
従ってselfの最高係数は
\[
 a_{12}=s,\quad a_{14}=\omega^2s,\quad a_7=\omega s,
 \qquad s\ne0
\]
で、他の $U^2c^2$ 係数は零である。
$L$ が固有vectorでないことから $m,n$ は共に非零。
対角generator/top変更で $s=m=n=T=1$ とできる。
例えば $c$ を固定し、$\gamma^3=Ts/(mn)$ として
$b,x,y$ をそれぞれ $\gamma,T/(\gamma m),T/(\gamma n)$ 倍、topを $1/(\gamma T)$ 倍すればよい。

元rank2条件の12線型式を係数体上で解くと、rankは7で、全八係数は一変数 $D$ により
\[
 (a_0,a_1,a_2,a_4,a_5,a_6,a_8,a_{15})
 =(D+\omega,D+\omega,-\omega(D+\omega),D,D,-\omega(D+\omega),\omega^2D,D)
\]
と書ける。
$A=D+\omega$ とすると元関係の像は、全 $D$ に対し
\begin{gather*}
 cx=\omega xc,\quad cy=-\omega yc,\quad
 (x+y)b=c^2,\quad b(x+y)=\omega^2c^2,\\
 bc=-(Ax+Dy)(x+y),\qquad
 cb=-\omega^2(x+y)(Ax+Dy)
\end{gather*}
である。
元画像を表示関係で掛け戻し、unit $6\times6$ minorとパラメータ分母がないことを確認する。

この一変数族の全次数3,4商から作る正常行列は一般点で $33\times7$ であり、
正常消去の7個のpivotは全て係数体の非零定数である。
全関係・正常pivotと分母の例外多項式は、平方因子を除いて
\[
 D(D+\omega/2)(D+\omega)
\]
となる。従ってこの三値以外では正常核は零。
三値を元商上で新しく構成すると、$D=0,-\omega$ は正常核零、
$D=-\omega/2$ は次数3商次元22で非AS4となる。
これで最後の全値を排除した。

公開再現器は \texttt{verify\_cubic\_third\_root\_nonunit\_zero\_column.py}。
元二列、自己対称graph、Cayley--Hamilton恒等式、全線型incidence、
元画像の全値同定、一変数全pivotと三つのfresh値を標数0で再生する。
通常のdomain・正常商の証明を、有限標本で置き換えてはいない。
\end{proof}

\begin{theorem}[三固有値H3包絡は存在しない]
\label{thm:cubic-three-eigenvalue-sector-empty}
強いintrinsic H3のAS4包絡の中山は、相異なる固有値をちょうど三つ持つことはない。
\end{theorem}
\begin{proof}
系\ref{cor:cubic-three-eigenvalue-order3-only}と
定理\ref{thm:cubic-third-root-unit-Jordan-empty}により、
この全領域の中山は半単純位数3、三つの必要スペクトルに限られる。
重複1は定理\ref{thm:cubic-third-root-unit-onezero-complete}で全排除済み。
非単位重複rootの元22tensor atlasは、$b^2$ 列の独立結合、非零比例結合、零結合に分かれる。
独立結合は定理\ref{thm:cubic-third-root-nonunit-independent-empty}、
非零比例結合は定理\ref{thm:cubic-third-root-nonunit-proportional-nondegenerate-empty}と
\ref{thm:cubic-third-root-nonunit-degenerate-self-empty}、
零結合は前定理により排除される。
一方だけ非零は純積関係で整域でない。
$\omega$ の共役にも同じ係数体上の証明が適用できるので、全三固有値領域が空である。
これは四単純・全CY・全成分と相対降下を完成したという主張ではない。
\end{proof}
